% !TEX program = pdflatex
% KAIROS --- Speaker notes / math companion
% Compile: pdflatex kairos_notes.tex (1x)

\documentclass[11pt]{article}

\usepackage[margin=1in]{geometry}
\usepackage{amsmath,amssymb,amsthm,mathtools}
\usepackage{xcolor}
\usepackage{hyperref}
\usepackage{enumitem}
\usepackage{booktabs}
\usepackage{lmodern}
\usepackage[T1]{fontenc}
\usepackage{microtype}
\usepackage{parskip}

\definecolor{KNavy}{HTML}{0B1E3F}
\definecolor{KGold}{HTML}{8A6B10}
\definecolor{KRule}{HTML}{6B5215}

\hypersetup{colorlinks=true,linkcolor=KNavy,urlcolor=KNavy}

\renewcommand{\section}[1]{\par\vspace{1.2em}\noindent{\color{KNavy}\bfseries\large #1}\par\vspace{0.2em}\noindent{\color{KRule}\rule{\textwidth}{0.4pt}}\par\vspace{0.4em}}
\renewcommand{\subsection}[1]{\par\vspace{0.8em}\noindent{\color{KNavy}\bfseries #1}\par\vspace{0.15em}}

\begin{document}

\noindent{\color{KNavy}\LARGE\bfseries KAIROS: math companion.}\\[0.3em]
\noindent{\color{KRule}\itshape Speaker notes for the colloquium equations. Not for the audience.}\\[0.3em]
\noindent{\color{KRule}\rule{\textwidth}{0.8pt}}

\vspace{0.5em}
\noindent A short walkthrough of every formula that appears in the deck,
with the intuition, the derivation, and the pitfalls I want to be ready
to defend if asked.

\section{1. Set-up: the position-by-time tensor}

For each gene $g$ we build a matrix
\[
    R_g \;\in\; \mathbb{R}^{M \times T},
    \qquad
    R_g[x,t] \;=\; \text{read count in bin } x \text{ at time } t.
\]

$M$ is the number of 250-bp position bins along the gene body (typically
50--300, depending on gene length). $T$ is the number of sampled time points
(here $T = 4$: 0, 10, 25, 40 minutes after estradiol).

The matrix is the full measurement. Everything downstream compresses it
into a summary statistic --- either position-by-position or gene-wide.

\section{2. The $1/v$ principle (the kinetic identity)}

The core physical intuition: at position $x$, a polymerase moving at speed
$v(x)$ spends time $1/v(x)$ per unit length. In a GRO-seq run-on snapshot,
the number of nascent RNA reads accumulated at that bin by time $t$ grows
proportionally to that dwell time, scaled by the initiation rate.

\[
    \mathbb{E}\!\left[R_g[x,t]\right] \;=\; \frac{c\, t}{v(x)}
    \qquad (c \text{ absorbs initiation rate and depth})
\]

Consequence: fit a straight line across the time course at each position:
\[
    R_g[x,t] \;=\; \alpha_g(x) \,+\, \beta_g(x)\,t \,+\, \varepsilon.
\]

The slope is the reciprocal of speed, up to the constant $c$:
\[
    \beta_g(x) \;=\; \frac{c}{v(x)}
    \qquad\Longrightarrow\qquad
    \widehat{v}_g(x) \;\propto\; \frac{1}{\widehat\beta_g(x)}.
\]

\subsection{Pitfall to flag.} The proportionality constant $c$ is not
recovered by the per-position regression. All $\widehat v(x)$ values are
velocities \emph{up to a gene-specific global scale}. This is why we
cannot compare absolute speeds across genes from $\widehat\beta$ alone.
We report \emph{relative} velocity profiles within a gene, and use
$\psi$ or DANKO for cross-gene magnitude.

\subsection{Pitfall to flag.} At early times and distal positions,
$R_g[x,t] = 0$ because the wave hasn't arrived. Fitting a slope through
an all-zero column is meaningless. We classify each position as
pre-wave, transition, or elongated (three-state logistic on cumulative
counts) and only fit slopes on transition-plus-elongated positions.

\section{3. Robust regression}

Because GRO-seq counts are heavy-tailed and PCR duplicates create
occasional spikes, we replace ordinary least squares with Huber's
M-estimator:
\[
    L_c(u) \;=\; \begin{cases}
        \tfrac12\,u^2, & \text{if } |u| \le c, \\
        c\,|u| - \tfrac12\,c^2, & \text{if } |u| > c,
    \end{cases}
\]

and solve $\widehat\beta(x) = \arg\min_\beta \sum_i L_c(R[x,t_i] -
\alpha - \beta t_i)$. Tuning $c = 1.345\,\widehat\sigma$ gives 95\%
efficiency at Gaussian noise while capping the leverage of a single
outlier at $c$.

\subsection{Why not Poisson GLM?} For $T=4$, the Poisson likelihood is
dominated by the small-count regime where maximum-likelihood estimates
are unstable; the Huber linear model is more forgiving and the slope
retains its read-per-minute interpretation.

\section{4. Gene-level summary I: algebraic diversity $Z_M$}

The per-position slope vector
$\widehat{\boldsymbol\beta}_g = (\widehat\beta_g(x_1),\ldots,\widehat\beta_g(x_M))$
lives in $\mathbb{R}^M$. We want one scalar per gene.

The diversity summary from Thornton \& Ray (2022) treats
$\{\widehat\beta_g(x)\}$ as elements of the abelian group
$(\mathbb{R},+)$ and counts effectively distinct group elements,
weighted by a continuous resolution function $f$:
\[
    Z_M \;=\; \frac{1}{M-1}\sum_{x=1}^{M}\mathbf{1}[\widehat\beta(x) \ne 0]\;
              f(\widehat\beta(x)).
\]

In practice we use $f(\beta) = |\beta|$ (magnitude) or $f(\beta) =
|\beta - \bar\beta|$ (spread around the mean). The estimator is
monotone in the number of positions with non-trivial signal and the
variation among them.

\subsection{Why not the mean slope?} The mean is dominated by pause spikes,
ignores concentration, and cannot distinguish one long slow patch from
many brief ones. $Z_M$ separates these.

\section{5. Gene-level summary II: spectral concentration $\psi$}

The more powerful summary comes from the singular-value decomposition of
the full tensor slice $R_g$:
\[
    R_g \;=\; \sum_{j=1}^{\min(M,T)} \lambda_j\,\mathbf{u}_j\,\mathbf{w}_j^{\!\top},
    \qquad \lambda_1 \ge \lambda_2 \ge \cdots \ge 0.
\]

\begin{itemize}[leftmargin=1.2em,itemsep=0.1em]
  \item $\mathbf{u}_j \in \mathbb{R}^M$ is the $j$-th \emph{position mode}.
  \item $\mathbf{w}_j \in \mathbb{R}^T$ is the $j$-th \emph{time mode}.
  \item $\lambda_j^2$ is the \emph{energy} of that mode in the total Frobenius norm.
\end{itemize}

The spectral concentration is the fraction of total energy in the first mode:
\[
    \psi(g) \;=\; \frac{\lambda_1^2}{\sum_j \lambda_j^2} \;\in\; (0,1].
\]

\subsection{Why this tracks wave-front speed.} A clean elongation wave
\emph{is} a rank-one matrix up to noise: the position profile $\mathbf{u}_1$
is the shape of the wave; the time profile $\mathbf{w}_1$ is its
growth through the time course. A gene with a single crisp wave has
$\psi \approx 1$. A gene where the wave is interrupted by multiple pauses,
or where read accumulation is diffuse across position and time, has
$\psi$ near the uniform value $1 / \min(M,T)$.

\subsection{Interpretation of the main result.} On 2{,}500 genes,
\[
    \mathrm{Spearman}(\psi,\ v^{\mathrm{DANKO}}) \;=\; 0.557,
    \qquad p < 10^{-12}.
\]
$\psi$ is not a substitute for the HMM's estimate of magnitude --- it
does not give you kilobases per minute. It gives you a \emph{ranking} of
genes by how coherently they elongate, and that ranking tracks a
standard wave-front rate remarkably well, for a fraction of the
compute.

\section{6. DANKO wave-front HMM (the anchor we validate against)}

For each gene, at each time point $t$, fit a two-state hidden Markov
model along position:
\begin{itemize}[leftmargin=1.2em,itemsep=0.1em]
  \item $S_0$ = \emph{pre-wave}, low emission mean $\mu_0$ (negative binomial).
  \item $S_1$ = \emph{post-wave}, high emission mean $\mu_1 \gg \mu_0$.
  \item Transition $S_0 \to S_1$ at probability $p_{01}$.
  \item $S_1$ is \emph{absorbing}: once elongation has started at a bin, it stays started.
\end{itemize}

Baum-Welch gives the maximum likelihood parameters; Viterbi gives the
most likely breakpoint position $\widehat x^\star_g(t)$. The gene-wide
elongation rate is simply
\[
    \widehat v_g^{\mathrm{DANKO}} \;=\; \widehat x^\star_g(t) / t,
\]
measured at $t = 40$ minutes.

\subsection{Why we reimplemented it in Python.} R groHMM is 14 minutes
single-threaded on 2{,}500 genes; our port is 7 minutes on 8 workers
with identical (Pearson $0.96$) output, and it exposes the internals
(forward-backward, emission model, Viterbi decoding) for extension.

\section{7. Epigenetic covariates: the mixed-effects model}

For each position $x$ in each gene $g$, assemble a feature vector
$\mathbf{z}_g(x) \in \mathbb{R}^K$ with $K = 10$ chromatin and sequence
covariates. The model:
\[
    \log \widehat v_g(x) \;=\; \mathbf{z}_g(x)^{\!\top}\boldsymbol\gamma
        \,+\, u_g \,+\, \eta(x) \,+\, \varepsilon.
\]

\begin{itemize}[leftmargin=1.2em,itemsep=0.1em]
  \item $\boldsymbol\gamma \in \mathbb{R}^{10}$: fixed effects of
        the $K$ covariates on $\log$ velocity.
  \item $u_g \sim \mathcal{N}(0,\sigma_g^2)$: gene-level random intercept.
  \item $\eta(x)$: positional spline soaking up TSS / polyA structure
        that is not covariate-driven.
  \item $\varepsilon \sim \mathcal{N}(0,\sigma^2)$: residual noise.
\end{itemize}

The log-link is important: it makes effects multiplicative and prevents
negative velocity estimates. Positive entries of $\widehat{\boldsymbol\gamma}$
are accelerators; negative entries are stallers.

\subsection{Top hits.} H3K36me3 ($+0.22$ per $z$), H3K27ac ($+0.14$),
DNase accessibility ($+0.11$) accelerate; H3K27me3 ($-0.18$), CpG
methylation ($-0.09$), GC content ($-0.08$) stall.

\section{8. Residual pause detection}

Pause candidates are positions where the observed velocity is much
lower than what the chromatin model predicts:
\[
    \rho_g(x) \;=\; \log \widehat v_g(x)
                 \,-\, \mathbf{z}_g(x)^{\!\top}\widehat{\boldsymbol\gamma}.
\]

Large-magnitude negative $\rho_g(x)$ flags a slow position that is
unexplained by local chromatin, i.e.\ a candidate regulated pause. We
call bins with $\rho < -1.8$ at FDR $< 0.05$.

\section{9. The recurring pattern}

Every KAIROS estimator follows the same recipe:

\begin{enumerate}[leftmargin=1.4em,itemsep=0.1em]
  \item Build the space-time tensor $R_g$ for each gene.
  \item Produce a position-indexed quantity from it: $\widehat\beta(x)$
        (slope), singular values $\lambda_j$, or HMM breakpoint.
  \item Compress to a scalar via linear algebra ($Z_M$, $\psi$, or $v/t$).
  \item Validate against an external anchor (DANKO HMM, simulation,
        chromatin covariates).
\end{enumerate}

The philosophy: replace hidden-state inference with spectral / algebraic
summaries wherever possible. Cheaper, more transparent, and in our
dataset, nearly as good.

\vfill
\noindent{\color{KRule}\rule{\textwidth}{0.4pt}}\\
{\color{KRule}\small Micah Thornton, PhD.\ \ Texas Woman's University.\ \
\texttt{github.com/mathornton01/ad-speed-profile-viewer}}

\end{document}
